% =====================================================================
%  Section 1 -- Notation and standing conventions
%  CT-S36-DRAFT-NOTATION.  Draft only.
%
%  Every symbol fixed here is fixed BEFORE its first use anywhere in the
%  paper.  The section exists because an isolated reader, given the
%  assembled material and nothing else, returned 171 items used before
%  being defined or never defined at all
%  (CT-S23-OUTSIDER-READ/outputs/coldread_main.txt, task 2).
% =====================================================================

\section{Notation and standing conventions}\label{sec:notation}

This section fixes every symbol, name and hypothesis used in the paper. It
introduces no mathematics; a reader who already knows the conventions may skip
to \S2 and return here as needed.

\subsection{The recurrence and its normalisation}\label{ssec:recurrence}

Throughout, $b$ is a polynomial with real coefficients,
\[
  b(n)=\sum_{k=0}^{d}\beta_k n^k ,\qquad \beta_d\neq0,\qquad d=\deg b\ge1 ,
\]
so that $\beta_d$ is the \emph{leading} coefficient and $\beta_{d-1}$ the
\emph{subleading} one. A \emph{polynomial continued fraction} (PCF) is the
continued fraction whose partial denominators are the values $b(1),b(2),\dots$
of such a $b$ at the positive integers; its denominators $Q_n$ satisfy the
three-term recurrence
\begin{equation}\label{eq:n-rec}
  Q_n=b(n)Q_{n-1}+Q_{n-2}\qquad(n\ge1),\qquad Q_{-1}=0,\quad Q_0=1 .
\end{equation}
Write
\begin{equation}\label{eq:n-norm}
  B_n=\prod_{k=1}^{n}b(k)\quad(B_0=1),\qquad
  \psi_n=\frac{Q_n}{B_n},\qquad
  \psi_\infty=\lim_{n\to\infty}\psi_n .
\end{equation}
We call $\psi_n$ the \emph{normalised denominator sequence} and $\psi_\infty$ its
\emph{limit}. The roots of $b$ are written
\begin{equation}\label{eq:n-roots}
  \nu_1,\dots,\nu_d\in\C,\qquad b(n)=\beta_d\prod_{i=1}^{d}(n-\nu_i).
\end{equation}

\paragraph{Three objects, one letter.}
The sequence \eqref{eq:n-norm}, its limit, and the roots \eqref{eq:n-roots} are
three different objects, and in the deposited record all three are written with
the letter $r$. That record's $r_1$ is therefore ambiguous: it is at once the
first term of the sequence and the first root.\footnote{%
  \textbf{Correspondence with the deposited record.} The names used in this
  paper and the names used in the Channel Theory deposit (Zenodo concept DOI
  \texttt{10.5281/zenodo.19941678}) correspond as follows:
  $\psi_n\leftrightarrow r_n$, $\psi_\infty\leftrightarrow r_\infty$, and
  $\nu_1,\dots,\nu_d\leftrightarrow r_1,\dots,r_d$. The deposited
  records are not altered; the renaming is made in this paper only, and this is
  the single place the correspondence is recorded. A fourth, unrelated object
  is written $r(\cdot)$ in \S\ref{ssec:katz} below, following van der
  Put--Saito; it is never one of these three.}
Because the collision is in the source material and not an artefact of
assembly, it is resolved here once and the paper uses $\psi_n$, $\psi_\infty$ and
$\nu_i$ everywhere.

\subsection{Standing hypotheses}\label{ssec:hypotheses}

Three hypotheses are used, and they are stated separately because they are
consumed separately. Only (A1) is in force by default; (A2) and (A3) are
invoked by name wherever they are needed.

\begin{itemize}
\item[\textbf{(A1)}] $b(k)\neq0$ for every integer $k\ge1$.
\item[\textbf{(A2)}] No root $\nu_i$ is a positive integer; equivalently,
  $\Gamma(1-\nu_i)$ is finite for every $i$. (The Gamma function has poles
  at $0,-1,-2,\dots$, so $\Gamma(1-\nu_i)$ is singular exactly when
  $\nu_i\in\{1,2,3,\dots\}$.) Note (A2) $\Rightarrow$ (A1).
\item[\textbf{(A3)}] $\psi_\infty=\lim_{n\to\infty}Q_n/B_n$ exists and is
  \emph{nonzero}.
\end{itemize}

(A2) is not decoration. It is exactly the condition under which
$\Gamma(1-\nu_i)$ is finite, and therefore the condition under which the
constant $K$ of \S\ref{ssec:growth} is defined at all. A statement that uses
$K$ and cites only (A1) has omitted a hypothesis.

\subsection{Asymptotic and analytic vocabulary}\label{ssec:asymptotic}

$\Gamma$ denotes Euler's Gamma function, the meromorphic continuation of
$\Gamma(s)=\int_0^\infty t^{s-1}e^{-t}\,dt$ ($\Re s>0$), with simple poles at
$s=0,-1,-2,\dots$ and $\Gamma(m+1)=m!$ for integers $m\ge0$. We write $(dn)!$
for $\Gamma(dn+1)$ when $dn$ is a non-negative integer.

For functions of $n\to\infty$: $f=O(g)$ means $|f|\le C|g|$ for some constant
$C$ and all large $n$; $f=o(g)$ means $f/g\to0$; and $f\sim g$ means
$f/g\to1$. In particular $o(1)$ denotes a quantity tending to $0$. These are
statements about $n\to\infty$ unless another limit is named.

A \emph{formal solution} is a solution taken as a formal power (or
power-and-exponential) series, with no convergence asserted. The
\emph{Newton polygon} of a linear differential operator at a point is the lower
convex hull, in the plane of (order of differentiation, order of vanishing) of
the coefficients, of the points attached to its terms; its slopes govern the
exponential growth of formal solutions at that point. A \emph{germ
coefficient} is a coefficient of such a local expansion, i.e.\ of the germ of
the solution at the point in question, as opposed to a global invariant. The
\emph{amplitude} of an exponentially small or large term is the prefactor
multiplying its exponential, and $c_\pm$ denote the amplitudes attached to the
two branches of a square-root (rank $\tfrac12$) irregular singularity. A
\emph{Stokes constant} is the coefficient with which one formal solution enters
the Borel-plane singularity structure of another across a Stokes line; it is a
connection datum, not a local one.

\subsection{The growth constants}\label{ssec:growth}

Under (A2) and (A3) the denominators satisfy
\[
  \frac{Q_n}{\Gamma(dn+1)}\;=\;K\cdot\Big(\frac{\beta_d}{d^{\,d}}\Big)^{n}
  \cdot n^{s}\cdot\bigl(1+O(1/n)\bigr) ,
\]
which defines the two constants used throughout: the \emph{growth exponent}
\[
  s=\frac{\beta_{d-1}}{\beta_d}+\frac{d-1}{2}
\]
and the \emph{connection constant} $K$, the amplitude of this asymptotic. $K$
depends on all of $\beta_0,\dots,\beta_d$; the base $\beta_d/d^{\,d}$ depends
only on $\beta_d$, and $s$ only on $\beta_d$ and $\beta_{d-1}$. Consequently a
test that inspects the base or the exponent alone is insensitive to
$\beta_0,\dots,\beta_{d-2}$, and cannot be used to detect their influence.

Two further symbols belong to the Borel plane of this asymptotic. We write
$\xi_0$ for the characteristic root
\[
  \xi_0=\frac{2}{\sqrt{\beta_2}}\qquad(d=2),
\]
and $\rho$ for the exponent of the leading Borel-plane singularity, an indicial
exponent in the uniformiser $u=\sqrt z$. An \emph{indicial exponent} of a
linear differential equation at a regular singular point is a root of its
indicial equation, equivalently an exponent $\alpha$ for which the equation
admits a local solution of the form $u^{\alpha}(c_0+c_1u+\cdots)$ with
$c_0\neq0$, possibly up to logarithmic terms. The letter $\rho$ is never used
for anything else; in particular it is not a root of $b$, which are the
$\nu_i$.

\subsection{Numerical vocabulary}\label{ssec:numerics}

\emph{dps} abbreviates \emph{decimal places of working precision}, in the sense
of \texttt{mpmath}'s \texttt{mp.dps}: a computation at dps $=60$ carries $60$
significant decimal digits throughout. \emph{Richardson extrapolation} is the
elimination of successive terms of an assumed expansion in powers of $1/n$ by
taking prescribed linear combinations of the values at several $n$; a
\emph{Richardson level} is one such elimination, so that $m$ levels remove the
terms in $n^{-1},\dots,n^{-m}$. \emph{PSLQ} is the integer-relation algorithm
of Ferguson and Bailey: given real numbers $y_1,\dots,y_m$ it searches for
integers $a_1,\dots,a_m$, not all zero, with $\sum a_i y_i=0$ to the working
precision. A PSLQ result is evidence for a relation and never a proof of one.

Where a numerical limit is reported, the \emph{$N/2$-versus-$N$ agreement
criterion} is the requirement that the value obtained from the first $N/2$
terms and the value obtained from all $N$ terms agree to a stated number of
digits; the reported number of digits is the number on which they agree, not
the working precision.

\subsection{The three objects, and their comparison}\label{ssec:objects}

Three objects are compared in this paper. They are named once, here.

\begin{itemize}
\item \textbf{$\mathsf A$} is the generating function $\sum_{n\ge0}Q_nz^n$ of
  the denominators \eqref{eq:n-rec}.
\item \textbf{$\mathsf B$} is the differential operator $L_\varphi$ attached to
  that generating function.
\item \textbf{$\mathsf C$} is the scalar second-order equation
  $y''+P_{\mathsf C}\,y'+Q_{\mathsf C}\,y=0$ obtained from it by the change of
  variable $x=s_++\delta z$, where $s_\pm=(-1\pm i\sqrt{11})/6$ and
  $\delta=-i\sqrt{11}/3$.
\end{itemize}

The letters $\mathsf A$, $\mathsf B$, $\mathsf C$ always denote these three
objects. They are set in sans serif to distinguish them from the coefficients
$A$, $A'$, $C$ of a scalar equation written $A\,W''+A'\,W'+C\,W=0$, which is a
different use of the same letters and occurs only in that displayed form.

Writing a scalar equation as $\partial^2+P\partial+Q$, its \emph{trace-free
normal form} is obtained by removing the $P$ term, and the resulting potential
\[
  R=\tfrac14P^2+\tfrac12P'-Q
\]
is the \emph{normal potential}. Two operators are \emph{projectively
equivalent} when their trace-free normal forms agree, equivalently when they
differ by a rank-one twist $y\mapsto fy$; the normal potential is blind to such
a twist, while the pair $(P,Q)$ is not.

With this vocabulary: $\mathsf B$ and $\mathsf C$ are \emph{projectively
equivalent under $z\mapsto1/x$}, their normal potentials being identically
equal. A statement that two objects are projectively equivalent is therefore
weaker than a statement that they are isomorphic, and the paper distinguishes
the two.

\subsection{Katz invariants and the classification table}\label{ssec:katz}

We follow van der Put and Saito, \emph{Moduli spaces for linear differential
equations and the Painlev\'e equations}, Ann.\ Inst.\ Fourier \textbf{59} (2009)
2611--2667, and use their notation unchanged.

For a differential module over the field of formal Laurent series at a point
$c$, the \emph{Katz invariant} $r(c)$ is the largest slope of the Newton
polygon at $c$: it is $0$ exactly when $c$ is regular or regular singular, and
otherwise is the positive rational measuring the irregularity there. The
notation is theirs, and the symbol $r$ in $r(c)$ is unrelated to the sequence
and limit of \S\ref{ssec:recurrence}, which is why those were renamed.

The \emph{Katz multiset} of an object is the collection of its Katz invariants
at the singular points, written in the order used by their Table~2.1, namely
\[
  \bigl(r(0),\,r(1),\,r(\infty),\,r(t)\bigr),
\]
with a dash in the $r(t)$ position when the singular set contains no fourth
point $t$. It is a multiset and not a tuple in the sense that the classification
depends on the values present, not on a labelling of the points. Table~2.1
attaches to each such multiset a Dynkin symbol and a Painlev\'e equation; the
two rows used in this paper are
\[
  \widetilde D_5:\ \mathrm{PV},\ (0,0,1,-),\ \dim P=3,
  \qquad
  \widetilde D_6:\ \mathrm{PV}_{\deg}=\mathrm{PIII}(D_6),\ (0,0,\tfrac12,-),\
  \dim P=2 .
\]
\emph{Painlev\'e V} (PV) is the fifth Painlev\'e equation, and $\widetilde D_5$
is the Dynkin symbol labelling its row; the two name the same row of Table~2.1
and are used interchangeably, as there. $\mathrm{PIII}(D_6)$, $\mathrm{PIII}(D_7)$
and $\mathrm{PIII}(D_8)$ denote the three degenerations of the third Painlev\'e
equation, again by their row labels.

\textbf{A matching multiset is a necessary condition only.} Table~2.1 assigns a
row to a multiset; agreement of multisets therefore \emph{excludes} every row
whose multiset differs, and \emph{assigns} none. No statement in this paper
infers a Painlev\'e label from a multiset match alone.

\subsection{Monodromy and trace coordinates}\label{ssec:traces}

On the monodromy side, $B_1$, $B_2$, $B_3$ are the monodromy matrices of the
local system, one for each of the three punctures, in van der Put--Saito's
notation at p.~2637. They carry $a_i=\operatorname{tr}(B_i)$, and the affine
cubic surface of monodromy data is coordinatised by
\begin{equation}\label{eq:n-traces}
  x_1=\operatorname{tr}(B_2B_3),\qquad
  x_2=\operatorname{tr}(B_1B_3),\qquad
  x_3=\operatorname{tr}(B_1B_2).
\end{equation}
These are the \emph{trace coordinates}, and \eqref{eq:n-traces} says what they
are traces of. Where the paper speaks of comparison quantities it means
$x_1,x_2,x_3$ and nothing else.

The letter $B$ is thus used for the monodromy matrices $B_i$ and, in
\eqref{eq:n-norm}, for the product $B_n$ of values of $b$. The two never occur
in the same argument; the monodromy matrices always carry an index in
$\{1,2,3\}$ and the product always an index $n$.

\subsection{What is not fixed here}\label{ssec:notfixed}

Two conventions are deliberately left to their points of use, because they are
local to one argument: the branch of $\zeta=w^{1/d}$, fixed where the Borel
variable is introduced, and the orientation of the loops generating the
monodromy, fixed where $M_\infty$ is defined. Both are stated there, and
neither is used before it is stated.
